\documentclass[10pt,a4paper]{amsart}
\usepackage{amsmath,amssymb,amsthm,graphicx}
% Retained preamble for the Stacks Project chapter {CH} (shared project preamble).
% The chapter file pinned-original-source.tex only carries \input{{preamble}} on its first line,
% so no preamble text lives inside the chapter itself.  The Stacks Project uses one single shared
% preamble for every chapter: the snapshot stages it three times (sources/stacks-chow/preamble.tex,
% sources/stacks-etale-cohomology/preamble.tex and focused-seeds/AGP/sources/stacks-snapshot-20261011/
% preamble.tex) and those copies are identical in size.  The block below is copied mechanically,
% line for line, from sources/stacks-chow/preamble.tex lines 14-18 and 154-217: the Xy-pic setup and
% the native theorem declarations together with the mathematical macros.  Nothing was added, removed
% or redefined; the wrapper class and the amsmath/amssymb/amsthm/graphicx base packages are supplied
% by the assembler.
% For commutative diagrams we use Xy-pic
% Chapter: stacks-spaces-more-morphisms.tex -- More on Morphisms of Spaces.
\usepackage[all]{xy}

% We use 2cell for 2-commutative diagrams.
\xyoption{2cell}

% Theorem environments.
%
\theoremstyle{plain}
\newtheorem{theorem}[subsection]{Theorem}
\newtheorem{proposition}[subsection]{Proposition}
\newtheorem{lemma}[subsection]{Lemma}

\theoremstyle{definition}
\newtheorem{definition}[subsection]{Definition}
\newtheorem{example}[subsection]{Example}
\newtheorem{exercise}[subsection]{Exercise}
\newtheorem{situation}[subsection]{Situation}

\theoremstyle{remark}
\newtheorem{remark}[subsection]{Remark}
\newtheorem{remarks}[subsection]{Remarks}

\numberwithin{equation}{subsection}

% Macros
%
\def\lim{\mathop{\mathrm{lim}}\nolimits}
\def\colim{\mathop{\mathrm{colim}}\nolimits}
\def\Spec{\mathop{\mathrm{Spec}}}
\def\Hom{\mathop{\mathrm{Hom}}\nolimits}
\def\Ext{\mathop{\mathrm{Ext}}\nolimits}
\def\SheafHom{\mathop{\mathcal{H}\!\mathit{om}}\nolimits}
\def\SheafExt{\mathop{\mathcal{E}\!\mathit{xt}}\nolimits}
\def\Sch{\mathit{Sch}}
\def\Mor{\mathop{\mathrm{Mor}}\nolimits}
\def\Ob{\mathop{\mathrm{Ob}}\nolimits}
\def\Sh{\mathop{\mathit{Sh}}\nolimits}
\def\NL{\mathop{N\!L}\nolimits}
\def\CH{\mathop{\mathrm{CH}}\nolimits}
\def\proetale{{pro\text{-}\acute{e}tale}}
\def\etale{{\acute{e}tale}}
\def\QCoh{\mathit{QCoh}}
\def\Ker{\mathop{\mathrm{Ker}}}
\def\Im{\mathop{\mathrm{Im}}}
\def\Coker{\mathop{\mathrm{Coker}}}
\def\Coim{\mathop{\mathrm{Coim}}}

% Boxtimes
%
\DeclareMathSymbol{\boxtimes}{\mathbin}{AMSa}{"02}

%
% Macros for moduli stacks/spaces
%
\def\QCohstack{\mathcal{QC}\!\mathit{oh}}
\def\Cohstack{\mathcal{C}\!\mathit{oh}}
\def\Spacesstack{\mathcal{S}\!\mathit{paces}}
\def\Quotfunctor{\mathrm{Quot}}
\def\Hilbfunctor{\mathrm{Hilb}}
\def\Curvesstack{\mathcal{C}\!\mathit{urves}}
\def\Polarizedstack{\mathcal{P}\!\mathit{olarized}}
\def\Complexesstack{\mathcal{C}\!\mathit{omplexes}}
% \Pic is the operator that assigns to X its picard group, usage \Pic(X)
% \Picardstack_{X/B} denotes the Picard stack of X over B
% \Picardfunctor_{X/B} denotes the Picard functor of X over B
\def\Pic{\mathop{\mathrm{Pic}}\nolimits}
\def\Picardstack{\mathcal{P}\!\mathit{ic}}
\def\Picardfunctor{\mathrm{Pic}}
\def\Deformationcategory{\mathcal{D}\!\mathit{ef}}
\usepackage{cleveref}
\crefname{originallemma}{original lemma}{original lemmas}
\Crefname{originallemma}{Original lemma}{Original lemmas}
\makeatletter
\newlabel{sites-modules-definition-morphism-locally-ringed-topoi}{{04HA}{source}{}{sourcepointer-1}{}}
\newlabel{sites-modules-definition-morphism-locally-ringed-topoi@cref}{{[originallemma][0][]04HA}{[1][0][]source}{}{}{}}
\makeatother
\crefname{originallemma}{original lemma}{original lemmas}
\Crefname{originallemma}{Original lemma}{Original lemmas}
\makeatletter
\newlabel{spaces-properties-lemma-morphism-locally-ringed}{{04KJ}{source}{}{sourcepointer-2}{}}
\newlabel{spaces-properties-lemma-morphism-locally-ringed@cref}{{[originallemma][0][]04KJ}{[1][0][]source}{}{}{}}
\makeatother
\crefname{originallemma}{original lemma}{original lemmas}
\Crefname{originallemma}{Original lemma}{Original lemmas}
\makeatletter
\newlabel{spaces-properties-theorem-fully-faithful}{{04KL}{source}{}{sourcepointer-3}{}}
\newlabel{spaces-properties-theorem-fully-faithful@cref}{{[originallemma][0][]04KL}{[1][0][]source}{}{}{}}
\makeatother
\title{A compatible map of structure sheaves of thickenings comes from a unique morphism of thickenings over the base}
\author{The Stacks Project Authors}
\date{stacks-project:spaces-more-morphisms.tex:local-snapshot-20261011; GNU FDL 1.2 or later}
\begin{document}
\maketitle
\noindent\textit{Editorial note.} All delivered mathematical text is the author's own, copied from the pinned Stacks Project chapter file spaces-more-morphisms.tex; no statement or proof step was rewritten, reordered or removed.\par\medskip
\noindent\textit{Editorial note.} Preamble provenance, as required for this source: the first line of the pinned chapter file is \textbackslash{}input\{preamble\}, and the native preamble was not staged with the source, so no preamble text lives inside the chapter. The retained preamble used here is the project's single Stacks-wide preamble (the same 217-line / 8116-byte file that the already published chow and etale-cohomology cards use), copied mechanically line for line from sources/stacks-chow/preamble.tex lines 14-18 (Xy-pic setup) and 154-217 (theorem declarations and mathematical macros). Not one character was changed, no package was added, no macro was defined, and the excerpt serves only as preamble, never as a span coordinate.\par\medskip
\noindent\textit{Editorial note.} The same-chapter pointers are located at their \textbackslash{}label lines inside the pinned chapter file; the cross-chapter pointers are located in the project's own tag list staged in the snapshot (sources/stacks-chow/tags.txt), at the verbatim lines carrying their canonical tags.\par\medskip
\noindent\textit{Editorial note.} The author leaves two verifications to the reader at the end of the proof: that f' o i\_X = i\_Y o f (obtained by a second use of the fully faithful theorem) and that f' is a morphism over B (proof lines 1642-1645). Both are routine checks of the construction that is carried out inside the span.\par\medskip
\noindent\textit{Allowed standard prerequisites (editorial).}\begin{itemize}
\item Morphisms between thickenings correspond to morphisms of the associated ringed topoi. Properties of Spaces, Theorem labelled spaces-properties-theorem-fully-faithful; invoked at lines 1624-1626 and again at 1643-1644. Reduces the existence and uniqueness of f' to producing a morphism of locally ringed topoi from the given algebra map alpha.
\item Definition of a morphism of locally ringed topoi. Modules on Sites, Definition labelled sites-modules-definition-morphism-locally-ringed-topoi; invoked at lines 1634-1636. Fixes exactly what has to be verified for (f\_\{spaces,etale\}, alpha) and uses the local nilpotence of the kernels of i\_X\textasciicircum{}sharp and i\_Y\textasciicircum{}sharp.
\item The underlying morphism is a morphism of locally ringed topoi. Properties of Spaces, Lemma labelled spaces-properties-lemma-morphism-locally-ringed; invoked at lines 1637-1639. Supplies the locally-ringed property for (f, f\textasciicircum{}sharp) against which alpha is compared.
\end{itemize}\medskip
\ifdefined\hypertarget\hypertarget{sourcepointer-1}{}\fi
\noindent\textit{Source reference.} Original lemma 04HA: sources/stacks-chow/tags.txt, source lines 5861--5861. Navigation pointer for \textbackslash{}ref\{sites-modules-definition-morphism-locally-ringed-topoi\} used inside the counted span (line 1636); the argument that consumes it is retained in the span.\par\medskip
\ifdefined\hypertarget\hypertarget{sourcepointer-2}{}\fi
\noindent\textit{Source reference.} Original lemma 04KJ: sources/stacks-chow/tags.txt, source lines 5975--5975. Navigation pointer for \textbackslash{}ref\{spaces-properties-lemma-morphism-locally-ringed\} used inside the counted span (line 1639); the argument that consumes it is retained in the span.\par\medskip
\ifdefined\hypertarget\hypertarget{sourcepointer-3}{}\fi
\noindent\textit{Source reference.} Original lemma 04KL: sources/stacks-chow/tags.txt, source lines 5977--5977. Navigation pointer for \textbackslash{}ref\{spaces-properties-theorem-fully-faithful\} used inside the counted span (line 1625); the argument that consumes it is retained in the span.\par\medskip
\begin{lemma}
\label{lemma-first-order-thickening-maps}
Let $S$ be a scheme. Let $B$ be an algebraic space over $S$.
Let $X \subset X'$ and $Y \subset Y'$ be thickenings
of algebraic spaces over $B$. Let $f : X \to Y$ be a morphism of algebraic
spaces over $B$. Given any map of $\mathcal{O}_B$-algebras
$$
\alpha : f_{spaces, \etale}^{-1}\mathcal{O}_{Y'} \to \mathcal{O}_{X'}
$$
such that
$$
\xymatrix{
f_{spaces, \etale}^{-1}\mathcal{O}_Y \ar[r]_-{f^\sharp} \ar[r] &
\mathcal{O}_X \\
f_{spaces, \etale}^{-1}\mathcal{O}_{Y'} \ar[r]^-\alpha
\ar[u]^{i_Y^\sharp} &
\mathcal{O}_{X'} \ar[u]_{i_X^\sharp}
}
$$
commutes, there exists a unique morphism of $(f, f')$ of
thickenings over $B$ such that $\alpha = (f')^\sharp$.
\end{lemma}

\begin{proof}
To find $f'$, by
Properties of Spaces, Theorem \ref{spaces-properties-theorem-fully-faithful},
all we have to do is show that the morphism of ringed topoi
$$
(f_{spaces, \etale}, \alpha) :
(\Sh(X_{spaces, \etale}), \mathcal{O}_{X'})
\longrightarrow
(\Sh(Y_{spaces, \etale}), \mathcal{O}_{Y'})
$$
is a morphism of locally ringed topoi. This follows directly
from the definition of morphisms of locally ringed topoi
(Modules on Sites,
Definition \ref{sites-modules-definition-morphism-locally-ringed-topoi}),
the fact that $(f, f^\sharp)$ is a morphism of locally ringed topoi
(Properties of Spaces,
Lemma \ref{spaces-properties-lemma-morphism-locally-ringed}),
that $\alpha$ fits into the given commutative diagram, and
the fact that the kernels of $i_X^\sharp$ and $i_Y^\sharp$ are
locally nilpotent. Finally, the fact that $f' \circ i_X = i_Y \circ f$
follows from the commutativity of the diagram and another application of
Properties of Spaces, Theorem \ref{spaces-properties-theorem-fully-faithful}.
We omit the verification that $f'$ is a morphism over $B$.
\end{proof}

\end{document}
